\chapter[Referencial Teórico]{Referencial Teórico}
\label{cap:referencial}

Este capítulo apresenta os conceitos que fundamentam a proposta. A
Seção~\ref{sec:plataformas} trata das plataformas de dados sobre as quais a
arquitetura se apoia. A Seção~\ref{sec:ml-producao} discute os desafios de
sistemas de ML em produção e os modelos de processo que organizam seu ciclo
de vida. A Seção~\ref{sec:mlops} define MLOps, seus princípios e componentes.
A Seção~\ref{sec:riscos} aborda os riscos específicos da modelagem com dados
longitudinais. A Seção~\ref{sec:setor-publico} discute o uso de ML no setor
público e a Seção~\ref{sec:relacionados} apresenta os trabalhos relacionados.

% ---------------------------------------------------------------------
\section{Plataformas de dados governamentais}
\label{sec:plataformas}

\subsection{Engenharia de dados e o ciclo de vida do dado}

\citeonline{reis2022} definem a engenharia de dados como o desenvolvimento,
a implantação e a manutenção de sistemas que recebem dados brutos e produzem
informação consistente para usos posteriores, como análise e aprendizado de
máquina. Os autores descrevem o ciclo de vida do dado em geração, ingestão,
armazenamento, transformação e disponibilização, permeado por correntes
transversais como segurança, gestão de dados, orquestração e boas práticas de
engenharia de software.

No contexto governamental, a etapa de ingestão é particularmente difícil. Os
dados vêm de sistemas construídos em épocas e por equipes diferentes, expostos
por APIs, arquivos, planilhas e extrações manuais, com vocabulários e
granularidades próprios. Fontes também mudam ao longo do tempo: no caso
estudado neste trabalho, a origem dos dados de situação das obras muda duas
vezes entre 2024 e 2026 (ver Seção~\ref{sec:caso}).

\subsection{Arquitetura em camadas e \textit{lakehouse}}

\citeonline{armbrust2021} propõem o \textit{lakehouse} como arquitetura que
combina o armazenamento aberto e barato de um \textit{data lake} com
funcionalidades de gestão típicas de um \textit{data warehouse}, como
transações, versionamento e controle de esquema. Os autores argumentam que
essa unificação atende, sobre os mesmos dados, tanto a análise via SQL quanto
cargas de trabalho de ciência de dados e ML.

Uma organização frequente nessas plataformas é a arquitetura em camadas
conhecida como medalhão, com três camadas de qualidade crescente:

\begin{itemize}
	\item \textbf{bronze}: preserva o dado bruto como recebido da fonte, sem
	modificações;
	\item \textbf{prata}: contém dados limpos, tipados e harmonizados entre
	fontes;
	\item \textbf{ouro}: reúne tabelas modeladas para consumo por análises,
	painéis e modelos.
\end{itemize}

A camada bronze é o ponto de reprocessamento: se uma regra de transformação
mudar, as demais camadas podem ser reconstruídas sem nova extração da fonte. O
Gov Hub adota essa organização, e a arquitetura proposta neste trabalho
(Capítulo~\ref{cap:proposta}) atribui a cada artefato do ciclo de vida de ML
uma dessas camadas.

\subsection{Orquestração e transformação declarativa}

Um orquestrador de fluxos de trabalho representa cada fluxo como um grafo
acíclico dirigido (DAG) de tarefas, com agendamento, novas tentativas e
registro de execuções. O Apache Airflow \cite{airflow} é um exemplo. Uma
ferramenta de transformação declarativa representa cada transformação como uma
consulta SQL versionada, resolve as dependências entre elas, materializa os
resultados e executa testes de dados declarados junto às transformações. O dbt
\cite{dbt} é um exemplo.

Os dois papéis separam duas responsabilidades: o orquestrador decide
\emph{quando} algo roda e em que ordem, e a camada de transformação define
\emph{o que} é calculado. Essa separação é relevante para ML, pois features
(variáveis de entrada de um modelo) são, em última análise, transformações
sobre dados. Catálogos de dados, como o OpenMetadata \cite{openmetadata},
completam essa organização registrando tabelas, linhagem e responsáveis.

% ---------------------------------------------------------------------
\section{Aprendizado de máquina em produção}
\label{sec:ml-producao}

\subsection{Dívida técnica em sistemas de ML}

\citeonline{sculley2015} mostram que, em sistemas reais, o código de ML é
uma pequena fração do sistema, cercada por coleta de dados, extração de
features, verificação de dados, gestão de configuração, infraestrutura de
serviço e monitoramento. Os autores identificam formas de dívida técnica
específicas desses sistemas, entre elas:

\begin{itemize}
	\item \textbf{Emaranhamento}: pelo princípio CACE (\textit{Changing Anything
	Changes Everything}), alterar uma feature, um hiperparâmetro ou a
	distribuição de entrada pode mudar o comportamento de todo o modelo;
	\item \textbf{Dependências de dados}: dependências instáveis ou
	subutilizadas entre o modelo e seus dados de entrada são mais difíceis de
	detectar do que dependências de código;
	\item \textbf{Código de cola e selvas de pipeline}: grande parte do sistema
	consiste em código que conecta componentes e prepara dados, frequentemente
	duplicado entre treino e uso em produção;
	\item \textbf{Ciclos de realimentação}: o próprio uso do modelo pode alterar
	os dados que ele observará no futuro.
\end{itemize}

\subsection{Desafios de implantação}

\citeonline{amershi2019}, em estudo de caso na Microsoft, descrevem um fluxo
de trabalho de ML em nove etapas e observam que gerenciar dados é mais difícil
do que gerenciar código, que a customização de modelos exige competências
diferentes das de engenharia de software tradicional e que os componentes de
ML se entrelaçam de formas que dificultam o isolamento.
\citeonline{paleyes2022} organizam os desafios relatados em estudos de caso
de implantação conforme as etapas de gestão de dados, aprendizado do modelo,
verificação e implantação, além de questões transversais como ética, aspectos
legais, confiança dos usuários finais e segurança. \citeonline{lwakatare2020}
chegam a uma classificação semelhante ao revisar sistemas de ML em larga
escala na indústria.

Em estudo de entrevistas com engenheiros de ML, \citeonline{shankar2024}
observam que a operação de modelos é iterativa, manual e coletiva, e que o
sucesso de uma implantação depende de três fatores: velocidade de
experimentação, validação antes da produção e versionamento de dados, código
e modelos.

\subsection{Modelos de processo}

O CRISP-DM (\textit{Cross-Industry Standard Process for Data Mining}) organiza
projetos de mineração de dados em seis fases iterativas: entendimento do
negócio, entendimento dos dados, preparação dos dados, modelagem, avaliação e
implantação \cite{wirth2000}. O CRISP-ML(Q) estende esse processo para ML,
acrescentando uma fase de monitoramento e manutenção e associando a cada fase
uma metodologia de garantia de qualidade baseada em riscos
\cite{studer2021}.

Neste trabalho, o CRISP-DM organiza o ciclo de vida de cada modelo, e as
fases de monitoramento e de garantia de qualidade do CRISP-ML(Q) orientam a
etapa de operação.

% ---------------------------------------------------------------------
\section{MLOps}
\label{sec:mlops}

\subsection{Definição e princípios}

\citeonline{kreuzberger2023}, a partir de revisão de literatura, análise de
ferramentas e entrevistas com especialistas, definem MLOps como um paradigma
que inclui boas práticas, conceitos e cultura de desenvolvimento voltados à
conceituação, implementação, monitoramento, implantação e escalabilidade de
produtos de ML de ponta a ponta. Os autores identificam nove princípios:

\begin{enumerate}
	\item[P1] automação de CI/CD;
	\item[P2] orquestração de fluxos de trabalho;
	\item[P3] reprodutibilidade;
	\item[P4] versionamento de dados, modelos e código;
	\item[P5] colaboração;
	\item[P6] treino e avaliação contínuos;
	\item[P7] rastreamento e registro de metadados de ML;
	\item[P8] monitoramento contínuo;
	\item[P9] ciclos de realimentação.
\end{enumerate}

Esses princípios se materializam em componentes técnicos: repositório de
código, componente de CI/CD, orquestrador de fluxos, repositório de features
(\textit{feature store}), infraestrutura de treino, registro de modelos,
repositório de metadados de ML, serviço de inferência e monitoramento
\cite{kreuzberger2023}. Os princípios P1--P9 servirão de critério de
avaliação da arquitetura proposta (Capítulo~\ref{cap:metodologia}).

\citeonline{testi2022} propõem uma taxonomia de trabalhos sobre MLOps e uma
metodologia para conduzir projetos de ML com práticas de MLOps. Em um dos
primeiros relatos brasileiros, \citeonline{aguiar2020} descreve a adoção de
MLOps por uma equipe com experiência prévia em DevOps durante 24 meses de
co-desenvolvimento entre governo e academia, e aponta que, embora implantar
modelos tenha ficado mais simples, manter e evoluir sistemas de ML segue
pouco discutido.

\subsection{Níveis de maturidade}

A documentação de arquitetura do Google Cloud descreve três níveis de
maturidade em MLOps \cite{google2020mlops}: no nível 0, o processo é manual
e o modelo é entregue como artefato; no nível 1, o pipeline de treino é
automatizado e permite treino contínuo; no nível 2, o próprio pipeline é
testado e implantado por CI/CD. \citeonline{john2021} propõem um modelo de
maturidade derivado de revisão de literatura e estudos de caso, com estágios
que vão de processos manuais orientados a dados até operação totalmente
automatizada com monitoramento e retreino.

O nível de maturidade adequado depende do contexto. Um modelo que gera
um escore mensal sobre dados atualizados mensalmente não exige inferência em
tempo real, mas exige que o retreino e o escore sejam reprodutíveis e
auditáveis.

\subsection{Rastreamento de experimentos e registro de modelos}

O MLflow \cite{zaharia2018} é uma plataforma aberta para o ciclo de vida de
ML com três componentes originais: \textit{Tracking}, que registra
parâmetros, métricas e artefatos de cada execução; \textit{Projects}, que
empacota código de forma reprodutível; e \textit{Models}, que define um
formato padrão para implantação. O componente de registro de modelos permite
versionar modelos e controlar sua promoção entre estágios.

\citeonline{breck2017} propõem o \textit{ML Test Score}, uma rubrica de 28
testes e verificações agrupados em features e dados, desenvolvimento do
modelo, infraestrutura de ML e monitoramento. A rubrica oferece um critério
objetivo para medir a prontidão de um sistema de ML para produção.

% ---------------------------------------------------------------------
\section{Riscos na modelagem com dados longitudinais}
\label{sec:riscos}

\subsection{Vazamento de informação}

\citeonline{kaufman2012} definem vazamento (\textit{leakage}) como a
introdução, nos dados de treino, de informação sobre o alvo que não estaria
legitimamente disponível no momento da predição. O vazamento produz modelos
que parecem excelentes na avaliação e falham em produção. Em dados
longitudinais, a forma mais comum é usar, para prever o estado no mês $t+h$,
features calculadas com informação posterior a $t$. A prevenção exige que
cada feature seja calculada ``na data'' (\textit{as of}) da observação, e que
a separação entre treino e teste respeite a ordem temporal.

\subsection{Validação temporal}

A validação cruzada com partições aleatórias supõe observações
independentes, hipótese que dados temporais violam. \citeonline{bergmeir2012}
comparam esquemas de avaliação para preditores de séries temporais, como a
validação cruzada em blocos e a avaliação fora da amostra sobre o trecho final
da série, e discutem em que condições cada um produz estimativas confiáveis.
Para o caso estudado, em que as features dependem do histórico de cada obra e
o objetivo é prever eventos futuros, adota-se a forma mais conservadora:
cortes temporais fixos e documentados, com treino sempre anterior à validação
e ao teste.

\subsection{Desalinhamento entre treino e produção}

\citeonline{polyzotis2018} descrevem o chamado \textit{training-serving
skew}, o desalinhamento entre treino e produção: diferenças entre os dados ou
o código usados para treinar o modelo e os usados para gerar predições. Uma
causa frequente é a reimplementação de uma mesma feature em dois lugares.
\citeonline{breck2019} propõem a validação sistemática de dados com esquemas
declarados e comparação de distribuições entre os conjuntos de treino e de
produção.

\subsection{Avaliação sob capacidade limitada e classes desbalanceadas}

Quando o modelo orienta uma ação de capacidade limitada, como vistorias, a
métrica relevante não é a acurácia global, e sim quantos eventos reais estão
entre os $k$ itens priorizados (\textit{recall@k}) ou qual a proporção de
acertos entre eles (\textit{precision@k}). Paralisações são eventos raros
em relação ao total de obras-mês, o que caracteriza um problema de classes
desbalanceadas. \citeonline{saito2015} mostram que, nesse cenário, a curva
precisão-revocação é mais informativa do que a curva ROC.

\subsection{Mudança de distribuição e monitoramento}

Modelos em produção enfrentam mudanças na distribuição dos dados de entrada
e na relação entre entradas e alvo (\textit{concept drift}).
\citeonline{gama2014} revisam técnicas de detecção e adaptação a essas
mudanças. \citeonline{rabanser2019} comparam empiricamente métodos de
detecção de mudança de distribuição e concluem que testes estatísticos de duas
amostras sobre representações de dimensão reduzida são uma linha de base
eficaz. No caso estudado, a troca da fonte de dados ao longo da série é uma
mudança de distribuição conhecida, que precisa ser tratada já na preparação
dos dados.

\subsection{Modelos para dados tabulares e explicabilidade}

Para dados tabulares, métodos de \textit{gradient boosting} sobre árvores de
decisão, como o LightGBM \cite{ke2017}, são uma escolha frequente pela boa
relação entre desempenho e custo computacional. Quando o interesse é o tempo
até um evento, a análise de sobrevivência oferece alternativas, do modelo de
riscos proporcionais de \citeonline{cox1972} a métodos de ML revisados por
\citeonline{wang2019survival}.

Para que um escore de risco seja usado por equipes de fiscalização, ele
precisa ser explicável. Os valores SHAP \cite{lundberg2017} atribuem a cada
feature uma contribuição para uma predição individual, com fundamentação na
teoria dos jogos cooperativos. \citeonline{mitchell2019} propõem os
\textit{model cards}, documentos curtos que acompanham um modelo e descrevem
seu uso pretendido, desempenho por subgrupos e limitações.

% ---------------------------------------------------------------------
\section{Aprendizado de máquina no setor público}
\label{sec:setor-publico}

\citeonline{wirtz2019} classificam as aplicações de inteligência artificial
no setor público e os desafios associados em quatro dimensões: implementação
tecnológica, lei e regulação, ética e sociedade. No Brasil, a Lei Geral de
Proteção de Dados \cite{brasil2018lgpd} impõe restrições ao tratamento de
dados pessoais, relevantes quando as bases de obras incluem informações de
construtoras e responsáveis técnicos.

\citeonline{ferreira2025}, ao analisar o pipeline de ML da plataforma Brasil
Participativo, concluem que ML responsável no setor público é sobretudo um
problema de engenharia institucional: acesso restrito a dados, conjuntos de
dados sem versão e governança frágil de proveniência e monitoramento pesam
tanto quanto as escolhas técnicas. Os autores defendem tratar pipelines de ML
como infraestrutura cívica, com práticas de dados transparentes,
reprodutíveis e auditáveis.

% ---------------------------------------------------------------------
\section{Trabalhos relacionados}
\label{sec:relacionados}

\citeonline{madaio2016} desenvolveram o Firebird, sistema que estima o risco
de incêndio de edificações comerciais em Atlanta para priorizar inspeções do
corpo de bombeiros. O problema é análogo ao deste trabalho: um escore de risco
orienta uma equipe de fiscalização de capacidade limitada. O trabalho, porém,
foca no modelo e na integração com o processo de inspeção, sem tratar da
arquitetura de operação contínua.

\citeonline{gondia2020} comparam árvores de decisão e classificadores
bayesianos ingênuos para prever o risco de atraso em projetos de construção a
partir de dados históricos de projetos. No contexto brasileiro,
\citeonline{fernandes2024} treinou modelos XGBoost sobre dados abertos do
SAGRES, do Tribunal de Contas do Estado da Paraíba, para classificar obras
públicas entre sucesso e fracasso. Ambos tratam a predição como experimento
isolado, sem pipeline de retreino, monitoramento ou integração com a
plataforma de dados de origem.

Do lado da engenharia, \citeonline{kreuzberger2023} propõem uma arquitetura
de referência genérica para MLOps, e \citeonline{aguiar2020} e
\citeonline{ferreira2025} relatam experiências de ML em projetos
governamentais brasileiros. Nenhum deles parte de uma plataforma de dados
governamentais já estruturada como \textit{lakehouse} em camadas medalhão para definir onde
cada responsabilidade do ciclo de vida de ML deve morar. A
Tabela~\ref{tab:relacionados} resume a comparação.

\begin{table}[htb]
	\centering
	\caption{Comparação entre trabalhos relacionados e este trabalho.}
	\label{tab:relacionados}
	\small
	\begin{tabularx}{\textwidth}{>{\raggedright\arraybackslash}p{3.8cm}*{4}{>{\centering\arraybackslash}X}}
		\toprule
		\textbf{Trabalho} & \textbf{Setor público} & \textbf{Predição de risco} & \textbf{Arquitetura MLOps} & \textbf{Plataforma de dados existente} \\
		\midrule
		\citeonline{madaio2016}      & Sim & Sim & Não     & Não \\
		\citeonline{gondia2020}      & Não & Sim & Não     & Não \\
		\citeonline{fernandes2024}   & Sim & Sim & Não     & Não \\
		\citeonline{kreuzberger2023} & Não & Não & Sim     & Não \\
		\citeonline{aguiar2020}      & Sim & Não & Parcial & Não \\
		\citeonline{ferreira2025}    & Sim & Não & Parcial & Não \\
		\textbf{Este trabalho}       & Sim & Sim & Sim     & Sim \\
		\bottomrule
	\end{tabularx}
	\fonte{Elaborado pelos autores.}
\end{table}

\pendente{Ampliar os trabalhos relacionados com uma busca estruturada
(Scopus, IEEE Xplore, ACM DL e SBC OpenLib) por ``MLOps'' AND (``public
sector'' OR ``government'') e por trabalhos de predição de paralisação ou
atraso de obras públicas no Brasil.}
